%% ----------------------------------------------------------------------------
%%  §10.3  La formula di Taylor con il resto integrale
%% ----------------------------------------------------------------------------
\section{La formula di Taylor con il resto integrale}
\label{sec:resto-integrale}

Il resto di Peano descrive il comportamento dell'errore $R_{n}(x)=f(x)-T_{n}(x)$ per
$x\to x_{0}$, ma non dice nulla sul suo valore per un $x$ \emph{fissato}: sapere che
$R_{n}(x)=o\bigl((x-x_{0})^{n}\bigr)$ non permette di affermare, ad esempio, che l'errore
commesso calcolando $\sen\frac1{10}$ con il polinomio $x-\frac{x^{3}}6$ è minore di $10^{-7}$. Per
questo occorre un'espressione esplicita del resto. Se la funzione è abbastanza regolare,
il calcolo integrale la fornisce.

\begin{teorema}[Formula di Taylor con il resto integrale]\label{teo:resto-integrale}
\index{formula di Taylor!con il resto integrale}
Sia $f$ derivabile $n+1$ volte nell'intervallo $[a,b]$, con derivata $f^{(n+1)}$ continua
(cioè $f$ di classe $C^{n+1}$ in $[a,b]$), e sia $x_{0}\in[a,b]$. Allora il resto
$R_{n}(x)=f(x)-T_{n}(x)$ della formula di Taylor di ordine $n$ con centro $x_{0}$ si
rappresenta nella forma
\begin{equation}\label{eq:resto-integrale}
  R_{n}(x)=\int_{x_{0}}^{x}\frac{(x-t)^{n}}{n!}\,f^{(n+1)}(t)\,dt \qquad\perogni x\in[a,b] .
\end{equation}
\end{teorema}

\begin{schema}
\Ipotesi{$f\in C^{n+1}([a,b])$; \ $x_{0}\in[a,b]$.}
\Tesi{$f(x)=\displaystyle\sum_{k=0}^{n}\frac{f^{(k)}(x_{0})}{k!}(x-x_{0})^{k}
+\int_{x_{0}}^{x}\frac{(x-t)^{n}}{n!}\,f^{(n+1)}(t)\,dt$ per ogni $x\in[a,b]$.}
\end{schema}

\begin{strategia}
Procediamo per induzione su $n$. Per $n=0$ la tesi è la formula fondamentale del calcolo
integrale: $f(x)-f(x_{0})=\int_{x_{0}}^{x}f'(t)\,dt$. Per passare da $n$ a $n+1$ osserviamo che
$R_{n}$ e $R_{n+1}$ differiscono per un solo termine,
$R_{n}(x)-R_{n+1}(x)=\frac{f^{(n+1)}(x_{0})}{(n+1)!}(x-x_{0})^{n+1}$: dobbiamo «estrarre» questo termine
dall'integrale \eqref{eq:resto-integrale}. Lo strumento è l'integrazione per parti nella
variabile $t$ (con $x$ fissato), derivando $f^{(n+1)}$ e integrando il peso
$\frac{(x-t)^{n}}{n!}$, la cui primitiva $-\frac{(x-t)^{n+1}}{(n+1)!}$ si annulla per $t=x$: in questo modo il
termine di bordo in $t=x$ scompare, e quello in $t=x_{0}$ è proprio il termine cercato.
\end{strategia}

\begin{dimostrazione}
Dimostriamo per induzione su $n\ge0$ la proposizione
\begin{quote}
$P(n)$: \emph{per ogni funzione $f$ di classe $C^{n+1}$ in $[a,b]$ e per ogni
$x_{0},x\in[a,b]$ vale la \eqref{eq:resto-integrale}.}
\end{quote}

\emph{Base} ($n=0$). Il polinomio di Taylor di ordine $0$ è la costante $f(x_{0})$, e
$R_{0}(x)=f(x)-f(x_{0})$. Poiché $f'$ è continua, per la formula fondamentale del calcolo
integrale (\cref{teo:formula-fondamentale})
\[
  \int_{x_{0}}^{x}f'(t)\,dt=\Bigl[f(t)\Bigr]_{t=x_{0}}^{t=x}=f(x)-f(x_{0})=R_{0}(x) ,
\]
che è la \eqref{eq:resto-integrale} per $n=0$ (si ricordi che $0!=1$ e $(x-t)^{0}=1$).

\emph{Passo induttivo.} Supponiamo vera $P(n)$ e sia $f$ di classe $C^{n+2}$ in $[a,b]$. In
particolare $f$ è di classe $C^{n+1}$, e per l'ipotesi induttiva
\[
  R_{n}(x)=\int_{x_{0}}^{x}\frac{(x-t)^{n}}{n!}\,f^{(n+1)}(t)\,dt .
\]
Fissato $x$, integriamo per parti (\cref{teo:per-parti-definito}) con
$F(t)=f^{(n+1)}(t)$, di classe $C^{1}$, e $G(t)=-\frac{(x-t)^{n+1}}{(n+1)!}$, che ha derivata
$G'(t)=\frac{(x-t)^{n}}{n!}$. Il teorema è enunciato per un intervallo $[\alpha,\beta]$ con $\alpha<\beta$, ma
vale anche quando l'estremo inferiore supera quello superiore, perché scambiando gli
estremi entrambi i membri cambiano segno; possiamo quindi usarlo anche se $x<x_{0}$.
Poiché $G(x)=0$,
\begin{align*}
  R_{n}(x)&=\int_{x_{0}}^{x}G'(t)F(t)\,dt
  =\Bigl[G(t)F(t)\Bigr]_{t=x_{0}}^{t=x}-\int_{x_{0}}^{x}G(t)F'(t)\,dt\\
  &=\frac{(x-x_{0})^{n+1}}{(n+1)!}\,f^{(n+1)}(x_{0})
  +\int_{x_{0}}^{x}\frac{(x-t)^{n+1}}{(n+1)!}\,f^{(n+2)}(t)\,dt .
\end{align*}
D'altra parte $T_{n+1}(x)=T_{n}(x)+\frac{f^{(n+1)}(x_{0})}{(n+1)!}(x-x_{0})^{n+1}$, e quindi
\[
\begin{aligned}
  R_{n+1}(x)&=R_{n}(x)-\frac{f^{(n+1)}(x_{0})}{(n+1)!}(x-x_{0})^{n+1}\\
  &=\int_{x_{0}}^{x}\frac{(x-t)^{n+1}}{(n+1)!}\,f^{(n+2)}(t)\,dt ,
\end{aligned}
\]
che è la \eqref{eq:resto-integrale} con $n+1$ al posto di $n$. Dunque vale $P(n+1)$, e per il
principio di induzione la tesi è dimostrata per ogni $n$.
\end{dimostrazione}

\begin{osservazione}[Il significato del peso]
La \eqref{eq:resto-integrale} dice che l'errore è una «media pesata» della derivata
$f^{(n+1)}$ lungo l'intervallo di estremi $x_{0}$ e $x$, con peso $\frac{(x-t)^{n}}{n!}$. Per $n=1$,
\[
  f(x)=f(x_{0})+f'(x_{0})(x-x_{0})+\int_{x_{0}}^{x}(x-t)\,f''(t)\,dt :
\]
la distanza tra il grafico e la retta tangente in $x_{0}$ dipende soltanto dalla derivata
seconda. Se $f''\ge0$, per $x>x_{0}$ l'integrando è non negativo, e per $x<x_{0}$ si ha
$\int_{x_{0}}^{x}(x-t)f''(t)\,dt=\int_{x}^{x_{0}}(t-x)f''(t)\,dt\ge0$: il grafico di una funzione con
derivata seconda non negativa sta sopra ciascuna delle sue rette tangenti, come sappiamo
dallo studio della convessità (\cref{cap:applicazioni-derivate}). Il peso $(x-t)$ è
massimo vicino a $x_{0}$ e nullo in $t=x$: la curvatura che si trova vicino al centro dello
sviluppo «pesa» di più, perché il suo effetto si accumula lungo tutto il tratto
successivo.
\end{osservazione}

\begin{esempio}[L'esponenziale supera i suoi polinomi di Taylor]\label{es:exp-maggiore-taylor}
Sia $f(x)=e^{x}$, $x_{0}=0$. Poiché $f^{(n+1)}(t)=e^{t}>0$, per $x>0$ l'integrando della
\eqref{eq:resto-integrale} è positivo in $(0,x)$, e quindi $R_{n}(x)>0$:
\begin{equation}\label{eq:exp-maggiore-taylor}
  e^{x}>1+x+\frac{x^{2}}2+\dots+\frac{x^{n}}{n!} \qquad\perogni x>0,\ \perogni n\ge0 .
\end{equation}
In particolare, con $x=1$, $e>\sum_{k=0}^{n}\frac1{k!}$ per ogni $n$: le somme $\sum_{k=0}^{n}\frac1{k!}$ sono
approssimazioni \emph{per difetto} di $e$. Nel \cref{sec:tabulazione} maggioreremo
l'errore $R_{n}(1)$, e otterremo così anche approssimazioni per eccesso.
\end{esempio}

\subsection{Un'applicazione: la freccia di un'ala a sbalzo}

Il resto integrale ha un'interpretazione fisica immediata nella teoria della trave. Nel
\cref{sec:applicazioni-integrali} abbiamo modellato la semiala come una trave incastrata
alla radice $y=0$ e libera all'estremità $y=s$, caricata dalla portanza distribuita
$l(y)$, e ne abbiamo calcolato il momento flettente
$M(y)=\int_{y}^{s}(\eta-y)\,l(\eta)\,d\eta$ \eqref{eq:taglio-momento}. Indichiamo con $v(y)$ lo
spostamento verticale (la \emph{freccia}) della linea d'asse nella stazione $y$. Nella
teoria di Eulero--Bernoulli la curvatura $\frac1\rho$ della linea d'asse deformata è
proporzionale al momento flettente, $\frac1\rho=\frac M{EI}$, dove $E$ è il modulo di
elasticità del materiale e $I$ il momento d'inerzia del second'ordine della sezione
rispetto all'asse neutro; supponiamo per ora $E$ e $I$ costanti. La curvatura del grafico
di $v$ vale $\frac{v''}{(1+v'^{2})^{3/2}}$ e, per piccole rotazioni ($\abs{v'}\ll1$), si riduce a $v''$.
Con $v$ positivo verso l'alto e $M>0$ per una portanza diretta verso l'alto (la
convenzione del \cref{sec:applicazioni-integrali}; molti testi usano la convenzione
opposta e scrivono $v''=-\frac M{EI}$), si ottiene l'equazione della linea elastica
\begin{equation}\label{eq:linea-elastica}
  v''(y)=\frac{M(y)}{EI} .
\end{equation}
L'incastro impone $v(0)=0$ e $v'(0)=0$: alla radice l'ala non si sposta e non ruota.

\begin{proposizione}[Freccia di una trave a sbalzo]\label{prop:freccia}
Sia $l$ continua in $[0,s]$, e sia $v$ di classe $C^{2}$ in $[0,s]$, soluzione della
\eqref{eq:linea-elastica} con $v(0)=v'(0)=0$. Allora
\begin{equation}\label{eq:freccia}
  v(y)=\frac1{EI}\int_{0}^{y}(y-t)\,M(t)\,dt \qquad\perogni y\in[0,s] .
\end{equation}
In particolare, per un carico uniforme $l(y)=l$ la freccia all'estremità vale
$v(s)=\frac{l\,s^{4}}{8EI}$.
\end{proposizione}

\begin{schema}
\Ipotesi{$l$ continua; \ $v\in C^{2}([0,s])$; \ $v''=\frac M{EI}$; \ $v(0)=v'(0)=0$.}
\Tesi{$v(y)=\frac1{EI}\int_{0}^{y}(y-t)M(t)\,dt$; \ con carico uniforme $v(s)=\frac{ls^{4}}{8EI}$.}
\end{schema}

\begin{strategia}
È la formula di Taylor di ordine $1$ con centro nell'incastro e resto integrale: il
polinomio di Taylor $v(0)+v'(0)y$ è nullo per le condizioni di incastro, e resta soltanto
il resto, in cui $v''$ si sostituisce con la \eqref{eq:linea-elastica}. Per il carico
uniforme si calcola $M$ e un integrale di potenze.
\end{strategia}

\begin{dimostrazione}
Per il \cref{teo:resto-integrale} con $n=1$ e $x_{0}=0$,
\[
  v(y)=v(0)+v'(0)\,y+\int_{0}^{y}(y-t)\,v''(t)\,dt=\frac1{EI}\int_{0}^{y}(y-t)\,M(t)\,dt .
\]
Se $l$ è costante, $M(t)=l\int_{t}^{s}(\eta-t)\,d\eta=\frac l2(s-t)^{2}$, e
\[
  v(s)=\frac l{2EI}\int_{0}^{s}(s-t)^{3}\,dt=\frac l{2EI}\Bigl[-\frac{(s-t)^{4}}4\Bigr]_{0}^{s}
  =\frac{l\,s^{4}}{8EI} . \qedhere
\]
\end{dimostrazione}

\begin{figure}[htb]
\centering
\begin{tikzpicture}[x=1.6cm, y=1.6cm]
  \fill[pattern=north east lines, pattern color=grigio] (-0.25,-0.9) rectangle (0,1.3);
  \draw[thick] (0,-0.9) -- (0,1.3);
  \draw[tratteggio] (0,0) -- (4.2,0);
  \draw[blu, very thick, domain=0:4, samples=60]
    plot (\x, {0.9*(\x/4)^2*(6-4*(\x/4)+(\x/4)^2)/3});
  \foreach \y in {0.25,0.75,...,3.75} {
    \pgfmathsetmacro{\v}{0.9*(\y/4)^2*(6-4*(\y/4)+(\y/4)^2)/3}
    \pgfmathsetmacro{\h}{0.6*sqrt(1-(\y/4)^2)}
    \draw[rame, -{Stealth[length=1.6mm]}] (\y,\v-0.06-\h) -- (\y,\v-0.06);
  }
  \node[rame, font=\scriptsize] at (1.0,-0.62) {$l(y)$};
  \draw[grigio, {Stealth[length=1.2mm]}-{Stealth[length=1.2mm]}] (3,0) -- (3,0.601);
  \node[font=\scriptsize] at (3,0.8) {$v(y)$};
  \draw[{Stealth[length=1.6mm]}-{Stealth[length=1.6mm]}] (4.15,0) -- (4.15,0.9)
    node[midway, right, font=\scriptsize] {$v(s)$};
  \draw[-{Stealth[length=1.8mm]}] (0,-1.05) -- (4.4,-1.05) node[right, font=\scriptsize] {$y$};
  \draw (4,-1.0) -- (4,-1.1) node[below, font=\scriptsize] {$s$};
  \draw (3,-1.0) -- (3,-1.1) node[below, font=\scriptsize] {$y$};
  \node[font=\scriptsize, below] at (0,-1.1) {$0$};
\end{tikzpicture}
\caption{La semiala come trave a sbalzo: incastrata alla radice, si inflette verso l'alto
sotto la portanza distribuita $l(y)$. Il polinomio di Taylor della freccia $v$ con
centro nell'incastro è nullo, e la freccia coincide con il resto integrale
\eqref{eq:freccia}. Gli spostamenti sono fortemente esagerati, e la deformata e il
profilo del carico sono soltanto qualitativi.}
\label{fig:linea-elastica}
\end{figure}

La \eqref{eq:freccia} si legge così: il tratto $[t,t+dt]$ della trave, dove il momento vale
$M(t)$, ruota di un angolo $\frac{M(t)}{EI}\,dt$, e questa rotazione, propagata con braccio $y-t$
fino alla stazione $y$, produce lo spostamento $(y-t)\frac{M(t)}{EI}\,dt$. Il peso $(y-t)$ del
resto integrale è proprio questo braccio. Per la semiala del
\cref{es:taglio-ellittico} ($W=10\,000$~N, $b=10$~m) e una rigidezza flessionale
$EI=2\cdot10^{6}$~N\,m$^{2}$ (un valore puramente indicativo), il carico uniforme
$l=\frac Wb=1000$~N/m dà una freccia all'estremità di circa $3{,}9$~cm; con la
distribuzione ellittica, che concentra la portanza verso la radice, la freccia scende a
circa $3{,}0$~cm, cioè si riduce di circa il $23\%$ (il rapporto esatto è calcolato
nell'\cref{es:freccia-ellittica}). Il programma calcola le due frecce integrando
numericamente. Si noti infine che la dimostrazione della \cref{prop:freccia} non usa la
costanza di $EI$: per un'ala reale, rastremata, la rigidezza $EI(t)$ diminuisce verso
l'estremità, e la stessa formula di Taylor dà $v(y)=\int_{0}^{y}(y-t)\frac{M(t)}{EI(t)}\,dt$, purché
$EI$ sia continua e positiva.

\begin{python}[Freccia all'estremità della semiala]
import math
from scipy.integrate import quad

W, b, EI = 10000.0, 10.0, 2.0e6            # N, m, N*m^2
s = b / 2
carichi = {
    "uniforme": lambda y: W / b,
    "ellittico": lambda y: 4 * W / (math.pi * b)
                           * math.sqrt(max(0.0, 1 - (y / s)**2)),
}
v_tip = {}
for nome, l in carichi.items():
    M = lambda t: quad(lambda e: (e - t) * l(e), t, s)[0]
    v_tip[nome] = quad(lambda t: (s - t) * M(t), 0, s)[0] / EI
    print(f"{nome:9s}:  v(s) = {100 * v_tip[nome]:.3f} cm")
print(f"formula l*s^4/(8EI): {100 * (W/b) * s**4 / (8*EI):.3f} cm")
print(f"rapporto ellittico/uniforme = "
      f"{v_tip['ellittico'] / v_tip['uniforme']:.4f}")
\end{python}
\begin{risultato}
uniforme :  v(s) = 3.906 cm
ellittico:  v(s) = 3.022 cm
formula l*s^4/(8EI): 3.906 cm
rapporto ellittico/uniforme = 0.7736
\end{risultato}
